\documentclass[../../main.tex]{subfiles}
\begin{document}

\begin{flushright}
\footnotesize\textit{【自動文字起こし・要確認】}
\end{flushright}

{\bf [解]}
$x\le0$の時，$f(x)=-\left(x+\dfrac{1}{2}\right)^2+\dfrac{1}{4}$，$x\ge0$の時，$f(x)=\left(x-\dfrac{1}{2}\right)^2-\dfrac{1}{4}$である．$x\le0$の時，$f'(x)=-2\left(x+\dfrac{1}{2}\right)$だから，$x=-1$の接線は
\begin{align}
y=x+1
\end{align}
であるから，$x\ge0$での交点は$x=1+\sqrt{2}$である．

\begin{figure}[htb]
\centering
\begin{tikzpicture}[scale=1.4, >=stealth]
  \draw[->] (-1.8,-0.3) -- (3.2,-0.3+0.5) node[right] {$x$};
  \draw[->] (0,-0.6) -- (0,3.2) node[above] {$y$};
  \draw[domain=-1.5:0, smooth, thick] plot (\x,{-(\x+0.5)*(\x+0.5)+0.25});
  \draw[domain=0:2.42, smooth, thick] plot (\x,{(\x-0.5)*(\x-0.5)-0.25});
  \draw[domain=-1.3:2.6, smooth] plot (\x,{\x+1});
  \node at (-1,-0.6) {$-1$};
  \node at (2.42,-0.6) {$1+\sqrt{2}$};
\end{tikzpicture}
\caption{$y=f(x)$と$x=-1$における接線$y=x+1$で囲まれる部分（斜線部）}
\end{figure}

あって，求める面積$S$として
\begin{align}
S&=\int_{-1}^0\left\{(x+1)+(x^2+x)\right\}dx+\int_0^{1+\sqrt{2}}\left\{(x+1)-(x^2-x)\right\}dx\\
&=\int_{-1}^0(x+1)^2\,dx+\int_0^{1+\sqrt{2}}(-x^2+2x+1)\,dx\\
&=\dfrac{1}{3}+\Bigl[-\dfrac{1}{3}x^3+x^2+x\Bigr]_0^{1+\sqrt{2}}\\
&=\dfrac{1}{3}+(1+\sqrt{2})+(1+\sqrt{2})^2-\dfrac{1}{3}(1+\sqrt{2})^3\\
&=\dfrac{6+4\sqrt{2}}{3}．
\end{align}

\newpage

\end{document}
